\section{Exact shared-target rigidity}
\label{sec:shared-target-rigidity}

We return to the literal canonical parent set
\(\cT_X\) of \cref{sec:literal-packet}.  The positivity of the orbit
coordinate and the signed canonical determinant together make target
collisions much more rigid than a generic equality of integers.

\begin{theorem}[Ordered composite-key uniqueness]
\label{thm:composite-uniqueness}
Let
\[
 t=(\alpha,\gamma,j),\qquad
 t'=(\alpha',\gamma',j')
\]
be canonical parent keys with
\begin{equation}
 \nu(t)=\nu(t'),\qquad
 x_t=x_{t'},\qquad
 y_t=y_{t'}.
\label{eq:ordered-collision-hyp}
\end{equation}
Then \(t=t'\).  Equivalently, the map
\begin{equation}
 t\longmapsto(\nu(t),x_t,y_t)
\label{eq:composite-key-map}
\end{equation}
is injective on the canonical parent archive.
\end{theorem}

\begin{proof}
The two target equalities imply
\[
 m_\alpha j=m_{\alpha'}j',
 \qquad
 m_\gamma j=m_{\gamma'}j'.
\]
The ordered target pairs have the same gcd, say \(c\).  Equality of
the normalized determinants therefore gives
\[
 m_\alpha-m_\gamma
 =
 m_{\alpha'}-m_{\gamma'}=:d.
\]
The off-diagonal mask makes \(d\ne0\).  Subtracting the two target
equalities gives
\[
 dj=dj',
\]
so \(j=j'\).  It follows that
\[
 m_\alpha=m_{\alpha'},\qquad
 m_\gamma=m_{\gamma'}.
\]
The physical row-slope uniqueness then gives
\(\alpha=\alpha'\) and \(\gamma=\gamma'\).
\end{proof}

\begin{remark}[Canonical-parent premise]
\label{rem:uniqueness-premise}
The theorem is a statement about mathematical parent keys.  Two
computational expansion rows may have the same parent; they must be
inverse-aggregated before \eqref{eq:composite-key-map} is tested.
Conversely, equality of the two complex numbers \(u_t=u_{t'}\)
without equality of their ordered targets is merely a coefficient
coincidence and is not a physical duplicate in the present sense.
\end{remark}

For the most inclusive target-sharing statement, allow equality
between either slot of one ordered target pair and either slot of the
other.

\begin{definition}[Any-slot shared-target graph]
\label{def:target-graph}
The graph \(\mathscr G_{\tar}\) has vertex set \(\cT_X\).  Distinct
vertices \(t,t'\) are joined when
\[
 \nu(t)=\nu(t')
 \quad\text{and}\quad
 \{x_t,y_t\}\cap\{x_{t'},y_{t'}\}\ne\varnothing.
\]
\end{definition}

\begin{lemma}[A distinct pair shares at most one target]
\label{lem:one-shared-target}
Two distinct vertices of \(\mathscr G_{\tar}\) satisfy exactly one
of the four possible target equalities
\[
 x_t=x_{t'},\quad x_t=y_{t'},\quad
 y_t=x_{t'},\quad y_t=y_{t'}.
\]
In particular, no off-diagonal edge is counted by two target-slot
relations.
\end{lemma}

\begin{proof}
Because \(j>0\) and \(m_\alpha\ne m_\gamma\), one has
\[
 x_t-y_t=(m_\alpha-m_\gamma)j\ne0,
\]
and similarly \(x_{t'}\ne y_{t'}\).  Thus two equalities that use the
same target slot of one vertex would force the two targets of the
other vertex to coincide.

It remains to consider two perfect matchings between the slots.  The
ordered matching
\[
 x_t=x_{t'},\qquad y_t=y_{t'}
\]
forces \(t=t'\) by \cref{thm:composite-uniqueness}.  For the crossed
matching
\[
 x_t=y_{t'},\qquad y_t=x_{t'},
\]
the contents agree.  Since the determinants agree, the signed row
differences agree:
\[
 m_\alpha-m_\gamma
 =
 m_{\alpha'}-m_{\gamma'}=:d\ne0.
\]
But target differences now give
\[
 dj=x_t-y_t
 =-(x_{t'}-y_{t'})
 =-dj'.
\]
Hence \(j'=-j\), contradicting \(j,j'>0\).
\end{proof}

\subsection{Divisor parametrization of one target group}

For \(s\in\{L,R\}\), write
\[
 T_L(t)=x_t,\qquad T_R(t)=y_t,
\]
and define
\begin{equation}
 G_s(k,T)
 =
 \{t\in\cT_X:\nu(t)=k,\ T_s(t)=T\},
\qquad
 g_s(k,T)=\#G_s(k,T).
\label{eq:target-groups}
\end{equation}

\begin{proposition}[Exact candidate parametrization]
\label{prop:target-parametrization}
Fix \(k\ne0\) and a positive target \(T\).

\begin{enumerate}[label=(\roman*),leftmargin=2.2em]
\item Every \(t=(\alpha,\gamma,j)\in G_L(k,T)\) is determined by a
  pair of positive divisors
  \[
   m_\alpha\mid T-h_0,\qquad c_t\mid T,
  \]
  through
  \[
   j=\frac{T-h_0}{m_\alpha},\qquad
   m_\gamma=m_\alpha-kc_t.
  \]
\item Every \(t\in G_R(k,T)\) is determined by
  \[
   m_\gamma\mid T-h_0,\qquad c_t\mid T,
  \]
  through
  \[
   j=\frac{T-h_0}{m_\gamma},\qquad
   m_\alpha=m_\gamma+kc_t.
  \]
\end{enumerate}
The formulas are necessary candidate conditions.  The physical row,
primitive, support, mask, and exact-gcd tests may only remove
candidates.
\end{proposition}

\begin{proof}
For the left group, \(T=m_\alpha j+h_0\), so
\(m_\alpha j=T-h_0\).  Also \(c_t=(x_t,y_t)\mid x_t=T\), and
\(\nu(t)=k\) gives
\[
 m_\alpha-m_\gamma=kc_t.
\]
The row slope determines the row key uniquely.  These equations prove
(i).  The right-hand argument is symmetric.
\end{proof}

\medskip
\begin{corollary}[Uniform divisor bound]
\label{cor:target-group-divisor-bound}
Uniformly in \(k\), \(T\), and the physical packet,
\begin{equation}
 \boxed{
 g_L(k,T),\ g_R(k,T)
 \le
 \tau(|T-h_0|)\tau(T)
 \ll_{\eps,h_0,\mathscr D}X^\eps.}
\label{eq:target-group-divisor-bound}
\end{equation}
\end{corollary}

\begin{proof}
The first estimate counts the divisor pairs in
\cref{prop:target-parametrization}.  On the physical support,
\(T\asymp_{\mathscr D}X\) and \(T-h_0\asymp_{\mathscr D}X\).
The standard divisor bound \(\tau(n)\ll_\eps n^\eps\) gives the last
estimate \citep{IwaniecKowalski2004}.
\end{proof}

\begin{remark}[What was not used]
No cancellation in \(u_t\), no distribution of primes, and no
average over \(h_0\) enters
\cref{prop:target-parametrization,cor:target-group-divisor-bound}.
The estimate is a structural count on one prescribed shift.  It is
also only an upper bound for a collision class; it does not lower-bound
the raw or post-bin mass.
\end{remark}
